% !TeX program = xelatex
% UTF-8; standalone XeLaTeX / XePersian transcription.
% Source: Google Drive / Notebook / sepsis &meningitis in neonates.
% pic01.jpg: 20261016_112148025.jpg — 1gcc0CcgKADrHNAsASrEi2oQzx33GsiLo
% pic02.jpg: 20261016_112155936.jpg — 1sFrq8HNfgGTLlDIuU7S7FBOur5KMAfQr
% pic03.jpg: 20261016_112159636.jpg — 1w_A346migDSid6UYX_pnsMwlep0K1AES
% pic04.jpg: 20261016_112204944.jpg — 1xSYoOuC9kqf06CJsJuLNGvPCCnShXvAB
% Values and statements are transcribed from the notes, without clinical correction.
% Dual-mode guard: compiles standalone with XeLaTeX, or is \input from the
% XePersian notebook (main.tex) where its local preamble below is skipped.
\newif\ifSepsisMeningitisStandalone
\ifdefined\settextfont
  \SepsisMeningitisStandalonefalse
\else
  \SepsisMeningitisStandalonetrue
\fi
\ifSepsisMeningitisStandalone
\documentclass[11pt,a4paper]{article}
\usepackage[margin=18mm]{geometry}
\usepackage{amsmath,amssymb}
\usepackage{enumitem}
\usepackage{xcolor}
\usepackage[normalem]{ulem}
\usepackage{xepersian}
\settextfont{Amiri}
\setlatintextfont{DejaVu Serif}
\setlength{\parindent}{0pt}
\setlength{\parskip}{4pt}
\setlength{\emergencystretch}{3em}
\setlist[itemize]{itemsep=4pt,topsep=4pt,leftmargin=*}
\newcommand{\en}[1]{\lr{#1}}
\newcommand{\marktext}[1]{\begingroup\setlength{\fboxsep}{1pt}\colorbox{yellow!40}{\strut #1}\endgroup}
\newcommand{\uncertain}[1]{\textcolor{red!65!black}{[خوانش نامطمئن: #1]}}
\newcommand{\Sepsissheet}[1]{%
  \section*{تصویر #1}
  {\small سربرگ: \en{Subject:} ۹۵ ـ سپسیس و مننژیت؛
  \en{Year:}\quad\en{Month:}\quad\en{Date:}؛
  شمارهٔ دست‌نویس پایین برگه: #1.\par}}

\begin{document}
\fi

% Helper macros used inside the body must exist in notebook mode too.
% (Each \providecommand must carry a full body: a body-less \providecommand
% consumes the following token as its definition and corrupts the document.)
\providecommand{\marktext}[1]{\begingroup\setlength{\fboxsep}{1pt}\colorbox{yellow!40}{\strut #1}\endgroup}
\providecommand{\uncertain}[1]{\textcolor{red!65!black}{[خوانش نامطمئن: #1]}}
\providecommand{\Sepsissheet}[1]{%
  \section*{تصویر #1}
  {\small سربرگ: \en{Subject:} ۹۵ ـ سپسیس و مننژیت؛
  \en{Year:}\quad\en{Month:}\quad\en{Date:}؛
  شمارهٔ دست‌نویس پایین برگه: #1.\par}}

\begin{center}
{\Large\bfseries ۹۵ ـ سپسیس و مننژیت}
\end{center}

\Sepsissheet{۱}
{\small یادداشت‌های کوچک کنار \en{Year} و \en{Month} در سربرگ تصویر اول:
\uncertain{۴۴؛ ۴۲؛ چند رقم و خط کوتاهِ ناخوانا}.\par}

\begin{itemize}
  \item ایمنی هومورال مادر از طریق \en{IgG}:
  \uline{استرپ گروه \en{B}}، \uline{\en{HSV}}.
  \item نوزادان با وزن پایین: خطر بالاتر ابتلا به بیماری‌ها به خاطر ضعف ایمنی؛
  ریسک ۶ برابری سپسیس در نوزادان زودرس.
  \item نوزادان زودرس زیر ۳۲ هفته \en{IgG} کافی ندارند؛
  \marktext{نیمهٔ دوم سه‌ماههٔ سوم}.
  \item نوزادان:
  \begin{itemize}
    \item ضعف در مهاجرت فاگوسیت‌ها به محل عفونت.
    \item ضعف در تولید اجزای کمپلمان.
    \item ضعف در ایمنی در مقابل ویروس‌ها.
  \end{itemize}
  \item \en{\(1/20\)} موارد سپسیس در نوزادان منجر به مننژیت می‌شود.
  \item شایع‌ترین ارگانیسم‌های عامل سپسیس و مننژیت در نوزادان:
  \begin{itemize}
    \item \marktext{استرپ گروه \en{B}}، \marktext{\en{E. coli}}، لیستریا مونوسایتوژنز.
    \item \marktext{ارگانیسم‌های گرم منفی}: کلبسیلا، سالمونلا، سراشیا.
  \end{itemize}
  \item در نوزادان با وزن تولد بسیار پایین:
  کوآگولاز منفی‌ها؛ \marktext{استافیلوکوک}.
  {\small (پیش از «نوزادان» یک واژه خط خورده است: \uncertain{در}.)}
\end{itemize}

\subsection*{سپسیس نوزادی}
\begin{itemize}
  \item \en{Early-onset sepsis}: \marktext{تولد تا روز ۷ام}.
  \item \en{Late onset sepsis}: \marktext{۸ تا ۲۸}.
  \item \en{Early-onset}: معمولاً در رحم؛ همراه با
  \marktext{عوامل باکتریال یا منشأ \en{GU}}.
  \item ارگانیسم‌های عامل سپسیس:
  \marktext{استرپ گروه \en{B}}، \en{E. coli}، کلبسیلا،
  ل. مونوسایتوژنز، \en{nt H. A.}
\end{itemize}

\clearpage
\Sepsissheet{۲}
\begin{itemize}
  \item اکثر نوزادان مبتلا به سپسیس \marktext{زودرس هستند}؛
  علائم غیر اختصاصی قلبی ـ ریوی: \marktext{تاکی‌پنه، سیانوز}.
\end{itemize}

\subsection*{ریسک فاکتورهای \en{Early onset sepsis}}
\begin{itemize}
  \item کلونیزیشن واژن با استرپ گروه \en{B}،
  \marktext{\en{PROM} بیش از ۲۴ ساعت}،
  آمنیونیت، لکوسیتوز یا تب در مادر، تاکی‌کاردی جنین، زایمان زودرس،
  جنسیت پسر و سیاه‌پوست بودن.
\end{itemize}

\subsection*{تظاهرات \en{Early onset sepsis}}
\begin{itemize}
  \item شایع‌ترین: \marktext{نارسایی تنفسی}، \marktext{شوک}، \marktext{مننژیت}.
  \item \marktext{\en{DIC}}، \marktext{\en{AKI}}،
  \marktext{گانگرن پریفرال قرینه}.
\end{itemize}

\subsection*{بررسی‌های مورد نیاز}
\begin{itemize}
  \item کشت خون و \en{CSF}، رنگ‌آمیزی گرم \en{CSF}، آنالیز \en{CSF}.
  \item \en{CSF} نرمال در نوزاد:
  \begin{itemize}
    \item پروتئین افزایش‌یافته: \en{100--150 mg/dl}.
    \item \en{WBC} افزایش‌یافته تا ۲۵ تا ۳۰؛ متوسط ۹ تا.
    \item ۷۵٪ لنفوسیت.
  \end{itemize}
  \item در مننژیت نوزادی:
  \marktext{نسبت گلوکز \en{CSF} به خون زیر ۵۰٪ است}.
  \item در این نوزادان باید \en{CBC} سریال اخذ شود؛ بررسی نوتروپنی:
  \begin{itemize}
    \item \marktext{باندسل} (نوتروفیل نابالغ).
    \item \marktext{ترومبوسایتوپنی}.
  \end{itemize}
  \item در سپسیس باکتریال \marktext{\en{CRP}} افزایش پیدا می‌کند.
  \item برای بررسی از نظر نومونی: \marktext{\en{CXR}}.
\end{itemize}

\clearpage
\Sepsissheet{۳}
\begin{itemize}
  \item در نوزاد با وزن بسیار پایین: خطر نومونی با پاتوژن‌های فلور نرمال واژن؛
  اوره‌آپلاسما اوره‌آلیتیکوم، مایکوپلاسما هومینیس.
  \item باید \en{ABG} سریال اخذ شود:
  \begin{itemize}
    \item \marktext{هایپوکسمی}.
    \item \marktext{اسیدوز متابولیک}.
    \item به علت: هایپوکسمی و شوک.
  \end{itemize}
\end{itemize}

\subsection*{درمان}
\begin{itemize}
  \item بخش اصلی درمان سپسیس و مننژیت:
  \marktext{درمان زودهنگام با \en{AB}}.
  \begin{itemize}
    \item \marktext{آمپی‌سیلین + آمینوگلیکوزید (جنتامایسین)}؛ برای ۱۰ تا ۱۴ روز.
    \item آلترناتیو: \marktext{آمپی‌سیلین + سفپیم}.
  \end{itemize}
  \item اگر مننژیت وجود داشت: تا ۲۱ روز یا ۱۴ روز پس از منفی شدن کشت \en{CSF}.
  \item مننژیت ناشی از استرپ گروه \en{B}:
  \begin{itemize}
    \item سن \en{\(\geq 7\)} روز:
    \en{high dose penicillin (300,000--450,000 U/kg/day)}،
    \en{divided every 6 hours}.
    \item سن \en{\(<7\)} روز: \en{375,000}.
  \end{itemize}
\end{itemize}

\subsection*{سپسیس تأخیری (۸ تا ۲۸)}
\begin{itemize}
  \item معمولاً در نوزادی که بدون مشکل ترخیص شده است.
  \item علائم بالینی: لتارژی، پورفیدینگ، هایپوتونی، آپنه، تشنج، تب،
  فونتانل بالج، هایپربیلی‌روبینمی دایرکت.
  \item \marktext{استئومیلیت}: \marktext{استاف اپیدرمیس، استرپ گروه \en{B}}.
\end{itemize}

\clearpage
\Sepsissheet{۴}
\begin{itemize}
  \item آرتریت: گونوکوک، استاف اورئوس، کاندیدا، گرم منفی‌ها.
  \item \en{UTI}: باکتری‌های \marktext{گرم منفی}.
  \item استئومیلیت ممکن است باعث سودوپارالیز شود.
  \item ارزیابی: \marktext{کشت ادرار} (استریل سوپراپوبیک آسپیراسیون)
  یا کاتتر ادراری.
  \item پاتوژن‌های \en{Late onset sepsis}:
  پاتوژن‌های \en{Early onset} + هموفیلوس آنفولانزا،
  استرپ نومونیه، نایسریا مننژیتیدیس؛ انتروویروس، \en{HSV}، \en{CMV}.
  \item درمان: \marktext{آمپی‌سیلین + سفالوسپورین نسل ۳}.
  \begin{itemize}
    \item در صورت وجود مننژیت: \marktext{ونکومایسین} هم اضافه می‌شود.
  \end{itemize}
\end{itemize}

\ifSepsisMeningitisStandalone
\end{document}
\fi
